\SourceLeaf{332}
datorum, subductis superfluis. esse debebit $=n$, ut ad coëfficientium
determinationem completam sufficiant. Iam quum ut superflui in hoc casu
reiiciendi sint valores functionis $X$ ii, qui respondent valori ipsius $x$
maiori quam $180^0$ nec non valori $0$ et $180^0$, habebimus pro quatuor
casibus supra distinctis:
\[
\begin{array}{rllcl}
1.&\text{quando }\mu\text{ par},&a=0,&\text{erit}&n=\frac12\mu-1\\
2.&\text{quando }\mu\text{ par},&a=\frac{180^0}{\mu},&\text{erit}&n=\frac12\mu\\
3.&\text{quando }\mu\text{ impar},&a=0.&\text{erit}&n=\frac12\mu-\frac12\\
4.&\text{quando }\mu\text{ impar},&a=\frac{180^0}{\mu},&\text{erit}&n=\frac12\mu-\frac12
\end{array}
\]
Iam prorsus simili modo ut in art. 37, in functione
\[
\begin{aligned}
&\gamma+\gamma'\cos x+\gamma''\cos 2x+\text{etc.}+\gamma^m\cos mx\\[6pt]
&{}+\delta'\sin x+\delta''\sin 2x+\text{etc.}+\delta^m\sin mx
\end{aligned}
\]
ad normam artt. 20, 22 eruta, quae omnibus $\mu$ valoribus satisfacit,
et in qua $m$ vel $=\frac12\mu-\frac12$, vel $=\frac12\mu$,
coëfficientes $\gamma$, $\gamma'$, $\gamma''$ etc. sponte evanescent.
Quum enim in serie $A,B,C,D$ etc. vel termini ultimi ordine retrogrado
in casu praesenti vel fiant $=-B,-C,-D$ etc. vel $=-A,-B,-C$ etc.,
prout $a=0$, vel $=\frac{180^0}{\mu}$, insuperque pro illo casu $A=0$,
manifesto
\[
\begin{aligned}
&A\cos\lambda a+B\cos\lambda b+C\cos\lambda c+D\cos\lambda d+\text{etc.}
\end{aligned}
\]
pro quovis valore ipsius $x$ erit $=0$. Quamobrem functio $X'=$
\[
\begin{aligned}
&\delta'\sin x+\delta''\sin 2x+\delta'''\sin 3x+\text{etc.}+\delta^m\sin mx
\end{aligned}
\]
in qua coëfficientes $\delta'$, $\delta''$, $\delta'''$ etc. determinantur
per aequationes
\[
\begin{aligned}
&\delta'=\frac2\mu(A\sin a+B\sin b+C\sin c+D\sin d+\text{etc.})\\[6pt]
&\delta''=\frac2\mu(A\sin 2a+B\sin 2b+C\sin 2c+D\sin 2d+\text{etc.})\\[6pt]
&\delta'''=\frac2\mu(A\sin 3a+B\sin 3b+C\sin 3c+D\sin 3d+\text{etc.})
\end{aligned}
\]
etc., ultimus autem. quando $\mu$ par est, adeoque $m=\frac12\mu$,
per hanc
\[
\begin{aligned}
&\delta^m=\frac1\mu(A\sin ma+B\sin mb+C\sin mc+D\sin md+\text{etc.})
\end{aligned}
\]
necessario cum $X$ identica erit, quod eodem modo, ut in art. 37,
facile demonstratur. Ceterum functio $X'$ in casu primo, ubi $a=0$,
$\mu$ par, revera ad or\SourceLeaf{333}%
dinem $\frac12\mu-1$ tantum ascendit, quum fiat
\[
\begin{aligned}
&\delta^m=\frac1\mu(A\sin 0+B\sin 180^0+C\sin 360^0+D\sin 540^0+\text{etc.})\\[6pt]
&=0
\end{aligned}
\]
\InterpolationArticle{40}
Si functio $X$ non est ordinis $n^{\mathrm{ti}}$, ut supposuimus,
sed ulterius excurrit, aequationes sequentes docebunt, quomodo differentia
inter $X'$ et $X$ a coëfficientibus sequentibus pendeat (v. artt. 21, 23)
\[
\begin{aligned}
&\delta'=\beta'\mp\beta^{\mu-1}\pm\beta^{\mu+1}-\beta^{2\mu-1}+\beta^{2\mu+1}\mp\text{etc.}\\[6pt]
&\delta''=\beta''\mp\beta^{\mu-2}\pm\beta^{\mu+2}-\beta^{2\mu-2}+\beta^{2\mu+2}\mp\text{etc.}\\[6pt]
&\delta'''=\beta'''\mp\beta^{\mu-3}\pm\beta^{\mu+3}-\beta^{2\mu-3}+\beta^{2\mu+3}\mp\text{etc.}
\end{aligned}
\]
et sic porro usque ad ultimum $\delta^m$ vel penultimum $\delta^{m-1}$,
prout $\mu$ impar est vel par; signa superiora hic valent, quando $a=0$,
inferiora, quando $a=\frac{180^0}{\mu}$: denique pro ultimo habetur
in casu (1) $\delta^m=0$, in casu (2) vero
$\delta^m=\beta^m-\beta^{3m}+\beta^{5m}-\beta^{7m}+\text{etc.}$
Omnes functiones periodicae, per quas $\mu$ valoribus propositis satisfit,
et ex quibus $X'$ est simplicissima, sub forma
$X'+Y\sin(\frac12\mu x-\frac12\mu a)$ sive generalius sub forma
$X''+Y\sin(\frac12\mu x-\frac12\mu a)$ contentae erunt, designante
$X''$ functionem talem quamcunque, quae formula pro casu 1 et 3 ad
$X''+Y\sin\frac12\mu x$, pro casu 2 et 4 autem ad
$X''+Y\cos\frac12\mu x$ reducitur; $Y$ vero, siquidem alias functiones
non consideramus, nisi quae ipsi $X$ sunt similes, i. e. e solis sinubus
compositae, necessario debet esse:
\[
\begin{array}{ll}
\text{in casu 1 formae}&g+g'\cos x+g''\cos 2x+\text{etc.}\\
\text{in casu 2 formae}&g\sin x+g'\sin 2x+g''\sin 3x+\text{etc.}\\
\text{in casu 3 formae}&g\cos\frac12x+g'\cos\frac32x+g''\cos\frac52x+\text{etc.}\\
\text{in casu 4 formae}&g\sin\frac12x+g'\sin\frac32x+g''\sin\frac52x+\text{etc.}
\end{array}
\]
Functionem $Y$ hic ita determinare licebit, ut prodeat functio ad ordinem
$\frac12\mu-1$, $\frac12\mu$, $\frac12\mu-\frac12$,
$\frac12\mu-\frac12$ depressa, quae cum $X'$ necessario identica erit.
Pro reductione functionis $X''$ ad $X'$ regula generalis sequens habetur:

Quivis terminus in $X''$ talis $L\sin\lambda x=L\sin(k\mu\pm\lambda')x$,
transmutetur aut in $\pm L\sin\lambda'x$ (quoties $a=0$, vel $k$ par),
aut in $\mp L\sin\lambda'x$ (quoties nec $a=0$, nec $k$ par,
i. e. quoties simul $a=\frac{180^0}{\mu}$ atque $k$ impar):
denique quoties in casu
\SourceLeaf{334}
primo i. e. pro \(a=0\), et valore pari ipsius \(\mu\) evadit \(\lambda'=\frac12\mu\), terminus \(L\sin\lambda x\) omnino destruatur.

\InterpolationArticle{41}

Quum omnes casus speciales in artt. 29---40 considerati ad casum generalem in artt. 20---28 absolute reducti sint, omnia artificia, per quae in hoc casu calculus abbreviatur, qualia in artt. 25, 26, 27 explicavimus, etiam ad illos applicari poterunt. Quamobrem non opus erit, huic disquisitioni immorari, cui sequens exemplum ad artt. 39, 40 pertinens, finem imponet.

Aequatio centri pro novo planeta \textit{Iunone}, adhibita excentricitate \(0{,}254236\), calculata est per methodum indirectam per singulos denos gradus, ut sequitur

\[
\begin{array}{r|r|rrrcr}
\multicolumn{2}{c|}{}&
\multicolumn{5}{c}{\text{Aequatio Centri}=X}\\
\multicolumn{2}{r|}{\text{Anomalia media}=x}&-&+&&&\\ \hline
0^{\mathord{0}}&360^{\mathord{0}}&0&&&&0\\
10&350&3^{\mathord{0}}&50'&38''30&\cdot\ \cdot\ \cdot\ \cdot&13838''30\\
20&340&7&38&21{,}47&\cdot\ \cdot\ \cdot\ \cdot&27501{,}47\\
30&330&11&20&8{,}79&\cdot\ \cdot\ \cdot\ \cdot&40808{,}79\\
40&320&14&52&48{,}06&\cdot\ \cdot\ \cdot\ \cdot&53568{,}06\\
50&310&18&12&49{,}21&\cdot\ \cdot\ \cdot\ \cdot&65569{,}21\\
60&300&21&16&17{,}02&\cdot\ \cdot\ \cdot\ \cdot&76577{,}02\\
70&290&23&58&42{,}92&\cdot\ \cdot\ \cdot\ \cdot&86322{,}92\\
80&280&26&14&55{,}85&\cdot\ \cdot\ \cdot\ \cdot&94495{,}85\\
90&270&27&58&52{,}36&\cdot\ \cdot\ \cdot\ \cdot&100732{,}36\\
100&260&29&3&28{,}13&\cdot\ \cdot\ \cdot\ \cdot&104608{,}13\\
110&250&29&20&33{,}68&\cdot\ \cdot\ \cdot\ \cdot&105633{,}68\\
120&240&28&41&2{,}10&\cdot\ \cdot\ \cdot\ \cdot&103262{,}10\\
130&230&26&55&22{,}77&\cdot\ \cdot\ \cdot\ \cdot&96922{,}77\\
140&220&23&55&2{,}70&\cdot\ \cdot\ \cdot\ \cdot&86102{,}70\\
150&210&19&35&0{,}79&\cdot\ \cdot\ \cdot\ \cdot&70500{,}79\\
160&200&13&57&40{,}52&\cdot\ \cdot\ \cdot\ \cdot&50260{,}52\\
170&190&7&16&58{,}33&\cdot\ \cdot\ \cdot\ \cdot&26218{,}33\\
180&180&0&&&&0
\end{array}
\]

Discerpimus hanc periodum \(36\) terminorum in sex minores senorum terminorum; valores seni functionis \(X\) in singulis periodis contenti exhibebuntur per formulam talem
\[
\begin{aligned}
&\gamma+\gamma'\cos x+\gamma''\cos 2x+\gamma'''\cos 3x\\[6pt]
&\quad+\delta'\sin x+\delta''\sin 2x+\delta'''\sin 3x
\end{aligned}
\]
\SourceLeaf{335}
ubi pro coëfficientibus \(\gamma,\gamma',\delta'\) etc. valores sequentes invenimus:


\begin{landscape}\thispagestyle{empty}\begingroup\fontsize{9}{12}\selectfont\setlength{\arraycolsep}{3pt}\renewcommand{\arraystretch}{1.3}
\[
\begin{array}{l|r|r|r|r|r|r|r|r}
\text{Periodus}&y=6x&\gamma&\gamma'&\gamma''&\gamma'''&\delta'&\delta''&\delta'''\\ \hline
\text{prima}&0^{\mathord{0}}&0&0&0&0&
-103830{,}165&
+15406{,}638&0\\
\text{secunda}&60&+56''205&-217''757&+761''671&-1528{,}825&-103937{,}346&+15841{,}387&-882{,}667\\
\text{tertia}&120&+56{,}115&-217{,}467&+760{,}805&-1527{,}397&-104151{,}314&+16709{,}402&-2645{,}530\\
\text{quarta}&180&0&0&0&0&-104258{,}100&+17142{,}684&-3525{,}740\\
\text{quinta}&240&-56{,}115&-217{,}467&-760{,}805&+1527{,}397&-104151{,}314&+16709{,}402&-2645{,}530\\
\text{sexta}&300&-56{,}205&-217{,}757&-761{,}671&+1528{,}825&-103937{,}346&+15841{,}387&-882{,}667
\end{array}
\]
\endgroup\end{landscape}


Singuli coëfficientes \(\gamma,\gamma',\gamma'',\gamma''',\delta',\delta'',\delta'''\) rursus sub formam talem
\[
\begin{aligned}
&\varepsilon+\varepsilon'\cos 6x+\varepsilon''\cos 12x+\varepsilon'''\cos 18x\\[6pt]
&\quad+\zeta'\sin 6x+\zeta''\sin 12x+\zeta'''\sin 18x
\end{aligned}
\]
reducentur: nullo vero negotio perspicitur, pro quatuor prioribus evanescere debere \(\varepsilon,\varepsilon',\varepsilon'',\varepsilon'''\); et pro tribus posterioribus, \(\zeta',\zeta'',\zeta'''\). Hoc modo invenitur
\[
\begin{aligned}
&\gamma=+64''{,}848\sin 6x+0''{,}052\sin 12x\\[6pt]
&\gamma'=-251''{,}277\sin 6x-0''{,}167\sin 12x\\[6pt]
&\gamma''=+879''{,}002\sin 6x+0''{,}500\sin 12x\\[6pt]
&\gamma'''=-1764''{,}511\sin 6x-0''{,}824\sin 12x\\[6pt]
&\delta'=-104044''{,}264+213''{,}968\cos 6x+0''{,}132\cos 12x\\[6pt]
&\qquad{}+0''{,}000\cos 18x\\[6pt]
&\delta''=+16275''{,}150-868''{,}020\cos 6x-0''{,}489\cos 12x\\[6pt]
&\qquad{}-0''{,}003\cos 18x\\[6pt]
&\delta'''=-1763''{,}689+1762''{,}868\cos 6x+0''{,}819\cos 12x\\[6pt]
&\qquad{}+0''{,}002\cos 18x
\end{aligned}
\]

His valoribus pro \(\gamma,\gamma'\) etc. substitutis, praeceptisque art. praec. observatis, prodit functio sequens pro aequatione centri, in qua singuli coëfficientes intra centesimam minuti secundi partem exacti sunt.


\[
\begin{array}{r|r}
-104044''264\sin x&-1''643\sin 9x\\
+16275{,}150\sin 2x&+0{,}494\sin 10x\\
-3527{,}378\sin 3x&-0{,}149\sin 11x\\
+873{,}511\sin 4x&+0{,}052\sin 12x\\
-232{,}622\sin 5x&-0{,}017\sin 13x\\
+64{,}848\sin 6x&+0{,}006\sin 14x\\
-18{,}655\sin 7x&-0{,}004\sin 15x\\
+5{,}491\sin 8x&+0{,}003\sin 16x
\end{array}
\]


\SourceLeaf{336}
Hoc modo manifesto calculus ita se habet, ac si pro arcu \(x\) alius \(x'=x-a\) introduceretur. Saepenumero etiam calculus magnopere sublevari potest, si pro functione \(X\) alia \(X-\mathfrak{x}\) adhibetur, ubi pro \(\mathfrak{x}\) functio quaelibet ipsi \(X\) similis assumi potest, modo ordinem \(\frac12\mu-\frac12{}^{\mathrm{tum}}\) pro valore impari vel ordinem \(\frac12\mu^{\mathrm{tum}}\) pro valore pari ipsius \(\mu\) non egrediatur, et in casu posteriori coëfficientes terminorum \(\cos\frac12\mu x\) et \(\sin\frac12\mu x\) (siquidem eo ascendit) teneant debitam rationem \((\cos\frac12\mu a:\sin\frac12\mu a)\). Manifestum est, si valores functionis \(\mathfrak{x}\) pro \(x=a.b,c,d\) etc. sint resp. \(\mathfrak{A},\mathfrak{B},\mathfrak{C},\mathfrak{D}\) etc., fore valores functionis \(X-\mathfrak{x}\) pro iisdem valoribus resp. \(A-\mathfrak{A},B-\mathfrak{B},C-\mathfrak{C},D-\mathfrak{D}\) etc.; si porro functio \(Z\) hisce valoribus satisfaciens per praecepta praecedentia investigatur, functio \(\mathfrak{x}+Z\) identica erit cum ea, quae ex applicatione directa horum praeceptorum ad valores ipsius functionis \(X\) prodiret. Semper autem functionem \(\mathfrak{x}\) ita determinare oportet, ut valores propositi functionis \(X\) prope per illam exhibeantur, quo pacto quantitates \(A-\mathfrak{A},B-\mathfrak{B},C-\mathfrak{C},D-\mathfrak{D}\) etc. parvae et ad calculum magis tractabiles evadent. Ceterum pro \(\mathfrak{x}\) etiam quantitas constans assumi potest, e. g. haec \(\frac1\mu(A+B+C+D+\text{etc.})\), quae functionis quaesitae partem invariabilem constituit.

\begin{center}\rule{0.28\linewidth}{0.4pt}\end{center}
\clearpage\fancyhead[L]{\small BEMERKUNGEN.}\pdfbookmark[0]{Bemerkungen zur Interpolation}{interpolation-bemerkungen}\SourceLeaf{337}
\begin{center}\textbf{BEMERKUNGEN.}\end{center}
\begin{center}\rule{3em}{0.3pt}\end{center}

Die vorliegende Abhandlung über Interpolation ist die wiederholte Ausarbeitung einer Untersuchung,
von welcher ein früherer Entwurf im October 1805 begonnen zu sein scheint. Sie enthält die Methoden, die
bei der Berechnung der Ausdrücke für die durch Störungskräfte hervorgebrachten Aenderungen der Elemente
einer Planetenbahn vielfach angewandt worden sind. Dieser Ausmittelung der Störungen steht ein anderes
Verfahren zur Seite, das direct die Werthe der Bahnelemente für einzelne Zeitpunkte ergiebt und das,
wie aus den im vorgefundenen handschriftlichen Nachlasse aufgezeichneten schliesslichen Formeln hervorgeht,
auf folgenden Principien beruht.

Bilden die Werthe des Arguments, für welche die Werthe der Function gegeben sind, eine arithmetische
Reihe, und hat man ein Argument $x$ eingeführt, das bei jenen Werthen wie die ganzen Zahlen
fortschreitet, bezeichnet die gegebenen Werthe mit $fx$, führt deren aufeinanderfolgende Differenzenreihen
ein:
\[
\begin{aligned}
&f(x+1)-fx=f^{1}(x+\frac12),\\[6pt]
&f^{1}(x+\frac12)-f^{1}(x-\frac12)=f^{2}x,\\[6pt]
&f^{2}(x+1)-f^{2}x=f^{3}(x+\frac12)\ \mbox{u.\,s.\,f.}
\end{aligned}
\]
ferner deren mit beliebigen Anfangsgliedern beginnende Summenreihen
$f^{-1}(x+\tfrac12),\ f^{-2}x,\ f^{-3}(x+\tfrac12)$ u.\,s.\,f.
für welche
\[
\begin{aligned}
&fx=f^{-1}(x+\frac12)-f^{-1}(x-\frac12),\\[6pt]
&f^{-1}(x+\frac12)=f^{-2}(x+1)-f^{-2}x\ \mbox{u.\,s.\,f.}
\end{aligned}
\]
endlich die Mittelwerthe der aufeinander folgenden Glieder der so entstandenen Reihen nemlich
\[
\begin{aligned}
&\mathinner{.\,.\,.} f^{-1}x=\frac12f^{-1}(x+\frac12)+\frac12f^{-1}(x-\frac12),\\[6pt]
&f(x+\frac12)=\frac12f(x+1)+\frac12fx,\\[6pt]
&f^{1}x=\frac12f^{1}(x+\frac12)+\frac12f^{1}(x-\frac12)\mathinner{.\,.\,.}
\end{aligned}
\]
also in übersichtlicher Anordnung zusammengestellt folgendes System von gegebenen und daraus zunächst
berechneten Werthen
\SourceLeaf{338}

\begin{landscape}\thispagestyle{empty}\begingroup\fontsize{9}{12}\selectfont\setlength{\arraycolsep}{3pt}\renewcommand{\arraystretch}{1.3}
\[
\begin{array}{ccc|cc||cc||cc|ccc}
 & \mbox{M. W.} & \mbox{II. Summe.} & \mbox{M. W.} & \mbox{I. Summe} & \mbox{gegeb. W.} & \mbox{M. W.} & \mbox{I. Differ.} & \mbox{M. W.} & \mbox{II. Differ.} & \mbox{M. W.} & \\
 & & \cdot & \cdot & & \cdot & & & \cdot & \cdot & & \\
 & \cdot & & & \cdot & & \cdot & \cdot & & & \cdot & \cdot \\
 & & f^{-2}(x-1) & f^{-1}(x-1) & & f(x-1) & & & f^{1}(x-1) & f^{2}(x-1) & & \\
\cdot & f^{-2}(x-\tfrac12) & & & f^{-1}(x-\tfrac12) & & f(x-\tfrac12) & f^{1}(x-\tfrac12) & & & f^{2}(x-\tfrac12) & \cdot \\
 & & f^{-2}x & f^{-1}x & & fx & & & f^{1}x & f^{2}x & & \\
\cdot & f^{-2}(x+\tfrac12) & & & f^{-1}(x+\tfrac12) & & f(x+\tfrac12) & f^{1}(x+\tfrac12) & & & f^{2}(x+\tfrac12) & \cdot \\
 & & f^{-2}(x+1) & f^{-1}(x+1) & & f(x+1) & & & f^{1}(x+1) & f^{2}(x+1) & & \\
\cdot & \cdot & & & \cdot & & \cdot & \cdot & & & \cdot & \cdot \\
 & & \cdot & \cdot & & \cdot & & & \cdot & \cdot & &
\end{array}
\]
\endgroup\end{landscape}

so gibt Art.\ 4, wenn man $1^{\mathrm{s}}$ für $t,\ a,\ b,\ c,\ d,\ e\ldotp\ldotp$
resp.\ $x+t,\ x,\ x+1,\ x-1,\ x+2,\ x-2$ $2^{\mathrm{s}}$ für
$t,\ a,\ b,\ c,\ d,\ e\ldotp\ldotp$ resp.\ $x+t,\ x,\ x-1,\ x+1,\ x-2,\ x+2\mathinner{.\,.\,.}$ setzt,
$3^{\mathrm{s}}$ aus den entsprechenden Seiten der beiden so erhaltenen Gleichungen die halbe Summe bildet,
$4^{\mathrm{s}}$ aus denselben Gleichungen, nachdem zuvor in der zweiten $x+1$ und $t-1$ resp.\ für $x$ und $t$ gesetzt ist,
die halbe Summe bildet und $5^{\mathrm{s}}$ in der zuletzt erhaltenen Gleichung $t=\tfrac12$ macht,
folgende Gleichungen für die durch Interpolation zu bestimmenden Werthe $\varphi(x+t)$ und
$\varphi(x+\tfrac12)$ derjenigen Function $\varphi x$, die bei jedem ganzzahligen Werthe
von $x$ der gegebenen Grösse $fx$ gleich wird:
\[
\begin{aligned}
&\varphi(x+t)={}fx+t\cdot f^{1}(x+\frac12)+\frac{t\cdot t-1}{1\cdot2}f^{2}x
\\[6pt]
&\qquad{}+\frac{t\cdot t-1\cdot t+1}{1\cdot2\cdot3}f^{3}(x+\frac12)+\ldotp\ldotp\\[6pt]
&+\frac{\Pi(t+n-1)}{\Pi 2n\cdot\Pi(t-n-1)}f^{2n}x
\\[6pt]
&\qquad{}+\frac{\Pi(t+n)}{\Pi(2n+1)\cdot\Pi(t-n-1)}f^{2n+1}(x+\frac12)\ldotp\ldotp
\end{aligned}
\]
\[
\begin{aligned}
&\varphi(x+t)={}fx+t\cdot f^{1}(x-\frac12)+\frac{t\cdot t+1}{1\cdot2}f^{2}x
\\[6pt]
&\qquad{}+\frac{t\cdot t+1\cdot t-1}{1\cdot2\cdot3}f^{3}(x-\frac12)+\ldotp\ldotp\\[6pt]
&+\frac{\Pi(t+n)}{\Pi 2n\cdot\Pi(t-n)}f^{2n}x
\\[6pt]
&\qquad{}+\frac{\Pi(t+n)}{\Pi(2n+1)\cdot\Pi(t-n-1)}f^{2n+1}(x-\frac12)\mathinner{.\,.\,.}
\end{aligned}
\]
\[
\begin{aligned}
&\varphi(x+t)={}fx+t\cdot f^{1}x+\frac{t\cdot t}{1\cdot2}f^{2}x
+\frac{t\cdot t-1\cdot t+1}{1\cdot2\cdot3}f^{3}x+\ldotp\ldotp\\[6pt]
&+\frac{\Pi(t+n-1)}{\Pi 2n\cdot\Pi(t-n)}t f^{2n}x
+\frac{\Pi(t+n)}{\Pi(2n+1)\Pi(t-n-1)}f^{2n+1}x\\[6pt]
&\qquad{}+\mathinner{.\,.\,.}
\end{aligned}
\]
\[
\begin{aligned}
&\varphi(x+t)={}f(x+\frac12)+(t-\frac12)f^{1}(x+\frac12)
+\frac{t\cdot t-1}{1\cdot2}f^{2}(x+\frac12)
\\[6pt]
&\qquad{}+\frac{t\cdot t-1}{1\cdot2}(t-\frac12)f^{3}(x+\frac12)+\ldotp\ldotp\\[6pt]
&+\frac{\Pi(t+n-1)}{\Pi 2n\cdot\Pi(t-n-1)}f^{2n}(x+\frac12)
\\[6pt]
&\qquad{}+\frac{\Pi(t+n-1)}{\Pi(2n+1)\cdot\Pi(t-n-1)}(t-\frac12)f^{2n+1}(x+\frac12)+\ldotp\ldotp
\end{aligned}
\]
\[
\begin{aligned}
&\varphi(x+\frac12)={}f(x+\frac12)-\frac18f^{2}(x+\frac12)
+\frac{3}{128}f^{4}(x+\frac12)\\[6pt]
&\qquad{}-\frac{5}{1024}f^{6}(x+\frac12)
+\frac{35}{32768}f^{8}(x+\frac12)\mathinner{.\,.\,.}\\[6pt]
&+(-1)^n\frac{\Pi 2n}{\Pi n\cdot\Pi n}\,2^{-4n}\cdot f^{2n}(x+\frac12)+\mathinner{.\,.\,.}
\end{aligned}
\]
Die Coëfficienten von $f^m$ in der dritten, vierten und fünften Gleichung sind gleich den Coëfficienten von
$(2i\sin\omega)^m$ in den Reihenentwicklungen resp.\ für
$\cos 2t\omega+i\dfrac{\sin 2t\omega}{\cos\omega},\allowbreak \quad
\dfrac{\cos(2t-1)\omega}{\cos\omega}+i\sin(2t-1)\omega$
und $\dfrac{1}{\cos\omega}$ nach Potenzen der Grösse $2i\sin\omega$.
\par\SourceLeaf{339}
Setzt man nun
\[
\begin{aligned}
&\Phi(x+\frac12)={}f^{-1}(x+\frac12)+\frac{1}{24}f^{1}(x+\frac12)
-\frac{17}{5760}f^{3}(x+\frac12)\\[6pt]
&\qquad{}+\frac{367}{967680}f^{5}(x+\frac12)
-\frac{27859}{164486400}f^{7}(x+\frac12)\\[6pt]
&+\frac{1295803}{122624409600}f^{9}(x+\frac12)-\ldotp\ldotp
\end{aligned}
\]
\[
\begin{aligned}
&\Phi x={}f^{-1}x-\frac{1}{12}f^{1}x+\frac{11}{720}f^{3}x
-\frac{191}{60480}f^{5}x+\frac{2497}{3628800}f^{7}x\\[6pt]
&-\frac{14797}{95800320}f^{9}x+\frac{92427157}{2615348736000}f^{11}x-\mathinner{.\,.\,.}
\end{aligned}
\]
\[
\begin{aligned}
&\Psi x={}f^{-2}x+\frac{1}{12}fx-\frac{1}{240}f^{2}x
+\frac{31}{60480}f^{4}x-\frac{289}{3628800}f^{6}x\\[6pt]
&+\frac{317}{22809600}f^{8}x-\ldotp\ldotp
\end{aligned}
\]
\[
\begin{aligned}
&\Psi(x+\frac12)=f^{-2}(x+\frac12)-\frac{1}{24}f(x+\frac12)
+\frac{17}{1920}f^{2}(x+\frac12)\\[6pt]
&\qquad{}-\frac{367}{193536}f^{4}(x+\frac12)+\ldotp-\ldotp\ldotp
\end{aligned}
\]
indem man als Coëfficienten von $f^m$ in diesen vier Gleichungen die Coëfficienten nimmt, die bei den
Entwicklungen resp.\ für
$\dfrac{1}{2i\omega},\allowbreak \ \dfrac{1}{2i\omega\cos\omega},\allowbreak \ -\dfrac{1}{4\omega\omega},\allowbreak \ -\dfrac{1}{4\omega\omega\cos\omega}$
noch Potenzen von $2i\sin\omega$ entstehen,
so sind $\Phi(x+\tfrac12),\ \Phi x,\ \Psi x,\ \Psi(x+\tfrac12)$ die besonderen Werthe solcher Functionen
$\Phi(x+t),\ \Psi(x+t)$,
für welche die Gleichungen
\[
\begin{aligned}
&\Phi(x+t)-\Phi(x+t_0)=\int_{t_0}^{t}\varphi(x+t)\,dt),\\[6pt]
&\Psi(x+t)-\Psi(x+t_1)=\int_{t_1}^{t}\Phi(x+t)\,dt
\end{aligned}
\]
Statt haben, sie werden also bestimmten Integralen von $\varphi(x+t)\,dt$ und
$\varphi(x+t)\cdot dt\cdot dt$ gleich, wenn
man die Anfangswerthe der Summenreihen
$f^{-1}(x-\tfrac12),\ f^{-1}(x+\tfrac12)\ldotp\ldotp,\ f^{-2}x,\ f^{-2}(x+1)\ldotp\ldotp$
auf geeignete Weise ausgewählt.

Die Ableitung der Reihen für $\Phi$ und $\Psi$ aus der dritten und vierten der obigen Gleichungen für
$\varphi(x+t)$ erhält man unmittelbar, wenn man die Integrationen so ausführt, dass zunächst die Ausdrücke für
$\Phi(x+\tfrac12)-\Phi(x-\tfrac12),\quad \Phi(x+1)-\Phi x,\quad
\Psi(x+1)-2\Psi x+\Psi(x-1)$ und
$\Psi(x+\tfrac32)-2\Psi(x+\tfrac12)+\Psi(x-\tfrac12)$
entstehen.
\begin{flushright}\textsc{Schering.}\end{flushright}
\begin{center}\rule{9em}{0.3pt}\end{center}
